\documentclass{article}
\usepackage[UTF8]{ctex}
\usepackage{amsmath,amssymb,amsfonts}
\usepackage{enumitem}
\usepackage{geometry}
\geometry{a4paper, margin=1in}
% === 错题本排版优化（不参与题目标识） ===
\setlength{\parindent}{0pt}
\setlength{\parskip}{0.35em plus 0.10em minus 0.08em}
\setlength{\emergencystretch}{3em}
\setlist{itemsep=0.25em, topsep=0.40em, parsep=0pt}
\allowdisplaybreaks[2]
% === 错题本排版优化结束 ===
\begin{document}
\section*{陕西师范大学2025年高等代数真题}
\section*{178．陕西师范大学2025年高等代数试卷}

\subsection{题目1}
一．（15 分）证明：对于任意的非负整数 \(n\) ，复数域上的多项式 \(x^{2}+x+1, x^{n+2}+(x+1)^{2 n+1}\) 都有

\[
x^{2}+x+1 \mid x^{n+2}+(x+1)^{2 n+1} .
\]

\subsection{题目2}
二．（15 分）计算行列式

\[
D_{n}=\left|\begin{array}{ccccc}
x & a & a & \cdots & a \\
-a & x & a & \cdots & a \\
-a & -a & x & \cdots & a \\
\vdots & \vdots & \vdots & & \vdots \\
-a & -a & -a & \cdots & x
\end{array}\right|
\]

\subsection{题目3}
三。（20分）设向量组 \(\alpha_{1}, \alpha_{2}, \cdots, \alpha_{n}(n \geq 3)\) 线性无关，且

\[
\beta_{1}=\alpha_{1}+\alpha_{2}, \beta_{2}=\alpha_{2}+\alpha_{3}, \cdots, \beta_{n-1}=\alpha_{n-1}+\alpha_{n}, \beta_{n}=\alpha_{1}+\alpha_{n} .
\]

讨论向量组 \(\beta_{1}, \beta_{2}, \cdots, \beta_{n}\) 的线性相关性。

\subsection{题目4}
四．（15 分）设 \(A\) 是秩为 \(m\) 的 \(m \times n\) 级矩阵，\(B\) 是秩为 \(n-m\) 的 \(n \times(n-m)\) 级矩阵，而且 \(A B=O\) 。若 \(n\) 维向量 \(\eta\) 是齐次线性方程组 \(A X=0\) 的解。证明：存在唯一的 \(n-m\) 维列向量 \(\xi\) ，使得 \(\eta=B \xi\) 。

\subsection{题目5}
五。（20 分）\(A, B\) 均是 \(4 \times 4\) 级实矩阵，齐次线性方程组 \(A X=0\) 的基础解系中包含3个解向量，\(B X=0\) 的一个基础解系中包含2个解向量。证明：

\begin{itemize}
  \item[（1）] \(A B X=0\) 至少有 3 个线性无关的解向量．
  \item[（2）] \(4 A+5 B\) 一定有实特征向量．
\end{itemize}

\subsection{题目6}
六。（15 分）设 \(A, B, C\) 是 \(n\) 级方阵，且 \(A\) 和 \(B\) 都可逆，证明：矩阵

\[
M=\left(\begin{array}{cc}
A & A \\
C-B & C
\end{array}\right)
\]

可逆，并求矩阵 \(M\) 的逆．

\subsection{题目7}
七．（15 分）\(V_{1}, V_{2}\) 是线性空间 \(V\) 的子空间，若 \(V_{1} \cup V_{2}\) 也是 \(V\) 的子空间．证明：

\begin{itemize}
  \item[（1）] \(V_{1} \cup V_{2}=V_{1}+V_{2}\) ．
  \item[（2）] \(V_{1} \subset V_{2}\) 或者 \(V_{2} \subset V_{1}\) ．
\end{itemize}

\subsection{题目8}
八．（20 分）设 \(V\) 是复数域上的一个 \(n\) 维线性空间，\(\sigma, \tau\) 为 \(V\) 上的两个线性变换，且 \(\sigma \tau=\tau \sigma\) ．证明：

\begin{itemize}
  \item[（1）] 若 \(\lambda_{0}\) 是 \(\sigma\) 的一个特征值，则特征子空间 \(V_{\lambda_{0}}\) 是 \(\tau\) 的一个不变子空间．
  \item[（2）] \(\sigma, \tau\) 至少有一个公共的特征向量．
\end{itemize}

\subsection{题目9}
九．（15 分）设 \(A, B\) 是 \(n\) 级正交方阵，若 \(|A| \neq|B|\) ，证明：\(A+B\) 是不可逆矩阵．
\end{document}
